\documentclass[11pt]{article}
\usepackage[margin=0.83in]{geometry}
\usepackage{amsmath,amssymb,amsthm,mathtools}
\usepackage{graphicx}
\usepackage{xcolor}
\usepackage[hidelinks]{hyperref}
\usepackage{microtype}

\newtheorem{theorem}{Theorem}
\newtheorem{lemma}{Lemma}
\newtheorem{corollary}{Corollary}
\newcommand{\B}{\mathbb B}
\newcommand{\R}{\mathbb R}
\newcommand{\Cc}{C_c^\infty}
\newcommand{\HH}{\mathcal H}
\newcommand{\kap}{\kappa_n}

\title{The higher-dimensional Hardy deficit is not a norm\\
\large A negative answer to the explicit norm subquestion in arXiv:1909.12587}
\author{Automated mathematical-discovery candidate (lane 11)}
\date{12 August 2026}

\begin{document}
\maketitle

\noindent\fcolorbox{black}{gray!8}{\parbox{0.94\linewidth}{
\textbf{Status.} Candidate counterexample; current verdict \textbf{likely valid}.
The result completely answers the paper's explicit norm subquestion for every
$n\ge3$, but it does \emph{not} settle existence of extremals for the singular
Hardy--Moser--Trudinger inequality.}}

\section{Source question and scope}

Nguyen's paper \cite{Nguyen2019} proves the sharp singular
Hardy--Moser--Trudinger inequality on $\B^n$.  On PDF page 5 it asks whether
extremals exist when $n\ge3$ and identifies a concrete obstruction: it was not
known whether $u\mapsto\HH(u)^{1/n}$ is a norm on $\Cc(\B^n)$, where
\begin{equation}\label{eq:H}
 \HH(u)=\int_{\B^n}|\nabla u|^n\,dx
 -c_n\int_{\B^n}\frac{|u|^n}{(1-|x|^2)^n}\,dx,
 \qquad c_n=\left(\frac{2(n-1)}n\right)^n .
\end{equation}
The same page proposes the analogous distance-to-the-boundary deficit on
convex domains.  The source passage is reproduced below.

\begin{center}
 \includegraphics[width=0.98\linewidth]{figures/open_problem_crop.png}
\end{center}

The result below answers the norm subquestion negatively.  It also shows that
the same obstruction persists for the source's natural smooth-convex-domain
functional.  The larger extremal-existence question remains open.

\section{Proof intuition}

If $N=\HH^{1/n}$ were a norm, then $\HH=N^n$ would be convex.  We disprove
convexity by looking at the second variation in a thin boundary annulus.  On
that annulus choose a base function which agrees with $1-|x|^2$.  After the
normal coordinate $s=1-|x|$ is introduced, its second variation is, up to a
positive asymptotically constant factor,
\[
 \int\left(|h'|^2-\kap\frac{|h|^2}{s^2}\right)ds,
 \qquad \kap=\left(\frac{n-1}{n}\right)^n.
\]
Smooth logarithmic boundary layers make the quotient of the two integrals
tend to the one-dimensional Hardy constant $1/4$.  But $\kap>1/4$ precisely
throughout the requested range $n\ge3$.  Thus the Hessian has a negative
direction, which is incompatible with any norm.  In a smooth domain, Fermi
coordinates give the same normal expression; curvature and tangential terms
are lower-order errors.

\section{A smooth logarithmic boundary layer}

\begin{lemma}\label{lem:layer}
For every $\delta>0$ there are nonzero $h_L\in\Cc(0,\delta)$ such that
\begin{equation}\label{eq:ratio}
 \frac{\int_0^\delta |h_L'(s)|^2\,ds}
 {\int_0^\delta |h_L(s)|^2s^{-2}\,ds}\longrightarrow\frac14.
\end{equation}
\end{lemma}

\begin{proof}
Fix nonzero $\chi\in\Cc(0,1)$ and put $\varepsilon_L=\delta e^{-L}$,
$g_L(t)=\chi(t/L)$, and
\[
 h_L(s)=s^{1/2}g_L\!\left(\log\frac{s}{\varepsilon_L}\right).
\]
Because $\chi$ vanishes in neighborhoods of both endpoints, $h_L$ is smooth
and compactly supported in $(\varepsilon_L,\delta)$.  With
$t=\log(s/\varepsilon_L)$,
\begin{align*}
 D_L&:=\int_0^\delta\frac{h_L^2}{s^2}\,ds
     =\int_0^L g_L^2\,dt=L\int_0^1\chi^2\,dt,\\
 \int_0^\delta |h_L'|^2\,ds
     &=\int_0^L\left(g_L'+\tfrac12g_L\right)^2dt
       =\tfrac14D_L+\int_0^L|g_L'|^2dt.
\end{align*}
The cross term vanishes by compact support, while
$\int|g_L'|^2=L^{-1}\int|\chi'|^2$.  Hence the quotient in
\eqref{eq:ratio} equals
\[
 \frac14+\frac1{L^2}
 \frac{\int_0^1|\chi'|^2}{\int_0^1|\chi|^2},
\]
which proves the claim.
\end{proof}

\section{Negative answer on the ball}

\begin{theorem}\label{thm:ball}
Let $n\ge3$.  The functional $\HH$ in \eqref{eq:H} is not convex on
$\Cc(\B^n)$.  Consequently $u\mapsto\HH(u)^{1/n}$ is not a norm on
$\Cc(\B^n)$.
\end{theorem}

\begin{proof}
Set $\kap=((n-1)/n)^n$.  The sequence $(1-1/n)^n$ is increasing for
$n\ge2$: for real $x>1$, the derivative of
$x\log(1-1/x)$ is $\log(1-1/x)+1/(x-1)>0$, by
$\log(1+y)<y$.  Hence
\begin{equation}\label{eq:gap}
 \kap\ge\left(\frac23\right)^3=\frac8{27}>\frac14.
\end{equation}

We first calculate with a radial direction $V(r)$ supported in a sufficiently
thin annulus next to $r=1$.  Choose $U\in\Cc(\B^n)$ radial and equal to
$1-r^2$ on a neighborhood of $\operatorname{supp}V$.  Such a $U$ exists
because each fixed support is a positive distance from the boundary.  On that
support, $U>0$ and $U'=-2r\ne0$.  Therefore differentiation under the integral
is legitimate and, writing $\omega_{n-1}=|S^{n-1}|$,
\begin{align}
 \frac{D^2\HH(U)[V,V]}{n(n-1)\omega_{n-1}}
 &=\int_0^1\left[(2r)^{n-2}(V')^2
 -c_n\frac{V^2}{(1-r^2)^2}\right]r^{n-1}\,dr.\label{eq:second}
\end{align}
Indeed, the scalar identity
$\frac{d^2}{dt^2}|a+tb|^n|_{t=0}=n(n-1)|a|^{n-2}b^2$
applies to both radial terms.

Put $s=1-r$ and $h(s)=V(1-s)$.  The right-hand side of
\eqref{eq:second} becomes
\begin{equation}\label{eq:qform}
 \int_0^1\left[A_n(s)|h'|^2-C_n(s)\frac{|h|^2}{s^2}\right]ds,
\end{equation}
where
\[
 A_n(s)=2^{n-2}(1-s)^{2n-3},\qquad
 C_n(s)=c_n\frac{(1-s)^{n-1}}{(2-s)^2}.
\]
As $s\downarrow0$,
\[
 A_n(s)\longrightarrow2^{n-2},\qquad
 C_n(s)\longrightarrow\frac{c_n}{4}=2^{n-2}\kap.
\]
Because of the strict gap in \eqref{eq:gap}, choose $\eta>0$ so small that
$(1+\eta)(1/4+\eta)<\kap-\eta$.  Then choose $\delta>0$ so that on
$(0,\delta)$,
\[
 A_n(s)\le2^{n-2}(1+\eta),\qquad
 C_n(s)\ge2^{n-2}(\kap-\eta).
\]
Finally take $h=h_L$ from Lemma~\ref{lem:layer} with $L$ large enough that
$\int|h'|^2<(1/4+\eta)\int h^2/s^2$.  Equation \eqref{eq:qform} is then
strictly negative.  Defining $V(r)=h(1-r)$ gives a member of
$\Cc(\B^n)$, and the compactly supported $U$ described above gives
\begin{equation}\label{eq:negative}
 D^2\HH(U)[V,V]<0.
\end{equation}

For sufficiently small $t>0$, Taylor's theorem and \eqref{eq:negative} give
\[
 \HH(U+tV)+\HH(U-tV)-2\HH(U)<0,
\]
which violates midpoint convexity.  If $N=\HH^{1/n}$ were a norm, then
\[
 \HH(\theta f+(1-\theta)g)=N(\theta f+(1-\theta)g)^n
 \le\bigl(\theta N(f)+(1-\theta)N(g)\bigr)^n
 \le\theta\HH(f)+(1-\theta)\HH(g),
\]
so $\HH=N^n$ would be convex.  This contradiction proves that $N$ is not a
norm.
\end{proof}

\section{Structure-preserving domain extension}

For a bounded smooth convex domain $\Omega\subset\R^n$, let
\begin{equation}\label{eq:domain}
 \HH_\Omega(u)=\int_\Omega|\nabla u|^n\,dx
 -\kap\int_\Omega\frac{|u|^n}{d(x,\partial\Omega)^n}\,dx.
\end{equation}
The strict Hardy inequality quoted in the source makes
$\HH_\Omega(u)^{1/n}$ pointwise well-defined on nonzero test functions.

\begin{corollary}\label{cor:domain}
For every bounded $C^\infty$ convex domain $\Omega\subset\R^n$ and every
$n\ge3$, $\HH_\Omega$ is nonconvex on $\Cc(\Omega)$.  Hence
$\HH_\Omega^{1/n}$ is not a norm.
\end{corollary}

\begin{proof}
Use Fermi coordinates $(y,s)$ in a smooth boundary patch, with
$s=d(x,\partial\Omega)$.  Choose $U\in\Cc(\Omega)$ which equals $s$ on a
neighborhood of the support of the direction below.  For
$V(y,s)=\zeta(y)h_L(s)$, where $\zeta$ is a fixed nonzero smooth tangential
cutoff, the same differentiation as above yields
\begin{equation}\label{eq:fermi}
 \frac{D^2\HH_\Omega(U)[V,V]}{n(n-1)}
 =\int\left[(\partial_sV)^2+\frac1{n-1}|\nabla_{\!T,s}V|^2
 -\kap\frac{V^2}{s^2}\right]J(y,s)\,dy\,ds.
\end{equation}
Here $J(y,s)=J(y,0)(1+O(s))$, and the tangential metrics are uniformly
equivalent for small $s$.  Let
$D_L=\int h_L^2/s^2$.  Lemma~\ref{lem:layer} gives
$\int|h_L'|^2=(1/4+o_L(1))D_L$, while
\[
 \int_0^\delta h_L^2\,ds\le\delta^2D_L.
\]
Consequently the Jacobian/metric errors in the normal and potential terms are
$O(\delta)D_L$, and the nonnegative tangential term is
$O(\delta^2)D_L$.  First choose $\delta$ so small that these errors are less
than half the strict gap $\kap-1/4$, and then choose $L$ large.  The expression
in \eqref{eq:fermi} is negative.  Thus $\HH_\Omega$ is nonconvex, and the
same norm-implies-convexity argument as in Theorem~\ref{thm:ball} finishes the
proof.
\end{proof}

\section{Verification, limitations, and novelty status}

The proof is analytic and has no computational dependency.  Verification
checked: (i) the exact second derivative factors; (ii) the limit
$c_n/(4\cdot2^{n-2})=\kap$; (iii) the strict threshold
$\kap>1/4$ for all $n\ge3$; (iv) genuine smooth compact support of every fixed
logarithmic layer and its localized base; and (v) the order of the Fermi-
coordinate errors.  No contradiction was found.

This packet does \emph{not} prove either existence or nonexistence of an
extremal for the singular inequality.  It shows instead that the particular
norm-completion route identified in the source cannot work: the proposed
gauge already fails the triangle inequality on test functions.  Likewise,
Corollary~\ref{cor:domain} does not decide the source's exponential inequality
on convex domains.

A bounded search on 12 August 2026 covered the exact title and arXiv id, the
quoted norm question, \texttt{Hardy--Moser--Trudinger extremals n>=3},
\texttt{Hardy deficit functional convexity p}, and
\texttt{Hardy functional does not define a norm}.  It found the source,
dimension-two results, and general Hardy papers, but no answer to this exact
higher-dimensional norm question and no matching nonconvexity theorem.
Novelty is therefore plausible, not certified.  Recommended human-review
focus is the compact-support localization in Theorem~\ref{thm:ball} and the
error bookkeeping in Corollary~\ref{cor:domain}.

\begin{thebibliography}{9}
\bibitem{Nguyen2019}
V. H. Nguyen,
\emph{The sharp Hardy--Moser--Trudinger inequality in dimension $n$},
arXiv:1909.12587 (2019), especially Theorem 1.1 and PDF page 5.
\end{thebibliography}

\end{document}
